\subsection{Pontryagin 对偶定理}

定理 2（Pontryagin）.　从 G 到 \(\widehat{G}\) 中的典范映射 η 是拓扑群的一个同构。它把 Haar 测度 dx 映为二重对偶 Haar 测度 d \(\widehat{x}\)。

我们先证明 \(\eta\) 是单射且是严格态射。为此只需证明：对 G 中 \(e\) 的每个邻域 U，存在 \(e\) 在 \(\widehat{\mathbf{G}}\) 中的一个邻域 W，使得 \(\overline{\eta}^{-1}(\mathrm{W})\subset \mathrm{U}\)（II，第 200 页，引理 2）。现设 V 是 G 中 \(e\) 的一个紧对称邻域，满足 \(\mathrm{V}^2\subset \mathrm{U}\)；设 \(f\) 是 G 上一个连续的正函数，其支集含于 V 中，且满足 \(f(e) > 0\)。设 \(g = \widetilde{f}*f\)。则 \(g\) 属于 \(\mathrm{A}(\mathrm{G})\)，其支集含于 U 中，且 \(g(e) > 0\)。此外，由 II，第 217 页，命题 11，\(\mathcal{F}_{\mathrm{G}}(g)\in \mathrm{L}^{1}(\widehat{\mathrm{G}})\)。集合 W 由这样的 \(\xi\)（属于 \(\widehat{\mathbf{G}}\)）组成，它使

\[
\left| \overline {{\mathcal {F}}} _ {\widehat {\mathrm{G}}} (\mathcal {F} _ {\mathrm{G}} (g)) (\xi) - \overline {{\mathcal {F}}} _ {\widehat {\mathrm{G}}} (\mathcal {F} _ {\mathrm{G}} (g)) (e) \right| <   \frac {1}{2} g (e)
\]

成立；W 是 \(e\) 在 \(\widehat{\mathbf{G}}\) 中的一个邻域，因为函数 \(\overline{\mathcal{F}}_{\widehat{\mathrm{G}}}(\mathcal{F}_{\mathrm{G}}(g))\) 在 \(\widehat{\mathbf{G}}\) 上连续。设 \(x \in \overline{\eta}^{-1}(\mathrm{W})\)。由公式 (26)，我们有

\[
\overline {{\mathcal {F}}} _ {\widehat {\mathrm{G}}} (\mathcal {F} _ {\mathrm{G}} (g)) (\eta (x)) = g (x),
\]

因而 \(|g(x) - g(e)| < \frac{1}{2} g(e)\)。这蕴含 \(g(x) \neq 0\)，从而 \(x \in \mathrm{U}\)，因为 \(g\) 的支集含于 U 中。于是 \(\overline{\eta}^{-1}(\mathrm{W}) \subset \mathrm{U}\)。

我们证明映射 \(\eta\) 是满射。由于该映射是到其像上的一个同胚，群 \(\eta(\mathbf{G})\) 是 \(\widehat{\widehat{\mathbf{G}}}\) 的一个局部紧子群，因此它在 \(\widehat{\widehat{\mathbf{G}}}\) 中是闭的（TG，III，第 22 页，推论 2）。我们用反证法，假设存在一个特征标 \(\xi \in \widehat{\widehat{\mathbf{G}}}\)，使得 \(\xi \notin \eta(\mathbf{G})\)。于是由 II，第 219 页，推论 1，存在一个非零元素 \(f\) 属于 \(\mathrm{L}^1(\widehat{\mathrm{G}})\)，使得 \(\mathcal{F}_{\widehat{\mathrm{G}}} (f)\) 在 \(\eta(\mathrm{G})\) 上为零。设 \(g \in \mathrm{L}^1(\mathrm{G})\)。函数 \((x, \chi) \mapsto g(x)f(\chi)\overline{\langle\chi,x\rangle}\) 属于 \(\mathrm{L}^1(\mathrm{G} \times \widehat{\mathrm{G}})\)。由 Lebesgue-Fubini 定理（INT，V，§8，第 4 节，定理 1, a)）可得

\[
\begin{array}{r} \int_ {\widehat {\mathrm{G}}} f (\chi) \mathcal {F} _ {\mathrm{G}} (g) (\chi) d \chi = \int_ {\mathrm{G}} g (x) \Bigl (\int_ {\widehat {\mathrm{G}}} f (\chi) \overline {{\langle \chi , x \rangle}} d \chi \Bigr) d x \\ = \int_ {\mathrm{G}} g (x) \mathcal {F} _ {\widehat {\mathrm{G}}} (f) (\eta (x)) d x = 0. \end{array}
\]

由于 Fourier 变换的像在 \(\mathcal{C}_{0}(\widehat{\mathrm{G}})\) 中稠密（II，第 209 页，命题 5 的推论），可知测度 \(f \cdot d\chi\) 为零。这与 \(f \neq 0\)（在 \(\mathrm{L}^{1}(\widehat{\mathrm{G}})\) 中）相矛盾，并证明了 \(\eta\) 是满射。

像测度 \(\eta(dx)\) 与测度 \(\nu\)（即测度 \(d\chi\) 的对偶测度）都是 \(\widehat{\widehat{G}}\) 上的 Haar 测度。设 \(f\) 是 A(G) 的一个非零元素；特别地 \(f \in \mathrm{L}^2(\mathrm{G})\)。由 II，第 219 页，命题 12，\(\mathcal{F}_{\mathrm{G}}(f) \in \mathrm{L}^{1}(\widehat{\mathrm{G}})\)，且

\[
\int_ {\widehat {\widehat {\mathrm{G}}}} \Big | \overline {{\mathcal {F}}} _ {\widehat {\mathrm{G}}} (\mathcal {F} _ {\mathrm{G}} (f)) \Big | ^ {2} \eta (d x) = \int_ {\mathrm{G}} | f | ^ {2} d x = \int_ {\widehat {\widehat {\mathrm{G}}}} \Big | \overline {{\mathcal {F}}} _ {\widehat {\mathrm{G}}} (\mathcal {F} _ {\mathrm{G}} (f)) \Big | ^ {2} d \nu ,
\]

其中第二个等式由两次应用 Plancherel 公式得到，因此 \(d\chi\) 的对偶 Haar 测度就是测度 \(\eta(dx)\)。

此后我们将把 G 与 \(\widehat{\mathrm{G}}\) 通过同构 \(\eta\) 等同起来。于是我们有：

推论.　从 \(L^{2}(\widehat{G})\) 到 \(L^{2}(G)\) 上的 Fourier 余变换与从 \(L^{2}(G)\) 到 \(L^{2}(\widehat{G})\) 上的 Fourier 变换互为逆等距。

注.　设 \(f \in \mathrm{L}^2(\mathrm{G})\)，\(g \in \mathrm{L}^2(\widehat{\mathrm{G}})\)。对 \(f\) 与 \(\overline{\mathcal{F}_{\widehat{\mathrm{G}}}(g)}\) 应用 Plancherel 公式 (24)，我们得到公式

\[
\int_ {\mathrm{G}} f (x) \mathcal {F} _ {\widehat {\mathrm{G}}} (g) (x) d x = \int_ {\widehat {\mathrm{G}}} \mathcal {F} _ {\mathrm{G}} (f) (\chi) g (\chi) d \widehat {x} (\chi)\tag{29}
\]

因为我们有 \(\mathcal{F}_{\mathrm{G}}(\overline{\mathcal{F}_{\widehat{\mathrm{G}}}(\boldsymbol{g})}) = \mathcal{F}_{\mathrm{G}}(\overline{\mathcal{F}}_{\widehat{\mathrm{G}}}(\overline{\boldsymbol{g}})) = \overline{\boldsymbol{g}}.\)

定义在 \(\mathcal{M}^{1}(\widehat{\mathrm{G}})\) 上的 Fourier 变换与余变换取值于 G 上连续有界函数所成的空间。对 \(\beta\in\mathcal{M}^{1}(\widehat{\mathrm{G}})\) 与 \(x\in\mathrm{G}\)，有

\[
\mathcal {F} _ {\widehat {\mathrm{G}}} (\beta) (x) = \int_ {\widehat {\mathrm{G}}} \overline {{\langle \chi , x \rangle}} d \beta (\chi), \qquad \overline {{\mathcal {F}}} _ {\widehat {\mathrm{G}}} (\beta) (x) = \int_ {\widehat {\mathrm{G}}} \langle \chi , x \rangle d \beta (\chi).
\]

G 与 \(\widehat{G}\) 的 Fourier 变换也互为转置。更精确地说：

命题 13.　设 \(\alpha \in \mathcal{M}^{1}(\mathrm{G})\)，\(\beta \in \mathcal{M}^{1}(\widehat{\mathrm{G}})\)。则我们有

\[
\mathcal {F} _ {\mathrm{G}} (\overline {{\mathcal {F}}} _ {\widehat {\mathrm{G}}} (\beta) \cdot \alpha) = \beta * \mathcal {F} _ {\mathrm{G}} (\alpha)\tag{30}
\]

特别地，有

\[
\int_ {\mathrm{G}} \mathcal {F} _ {\widehat {\mathrm{G}}} (\beta) (x) d \alpha (x) = \int_ {\widehat {\mathrm{G}}} \mathcal {F} _ {\mathrm{G}} (\alpha) (\chi) d \beta (\chi).\tag{31}
\]

在恒等式两边于 \(\chi = 1\) 处取值，即由公式 (30) 得出公式 (31)。我们来证明 (30)。设 \(\chi \in \widehat{G}\)。于是有

\[
\begin{array}{r l} & {(\mathcal {F} _ {\mathrm{G}} (\overline {{\mathcal {F}}} _ {\widehat {\mathrm{G}}} (\beta) \cdot \alpha)) (\chi) = \int_ {\mathrm{G}} \overline {{\langle \chi , x \rangle}} \overline {{\mathcal {F}}} _ {\widehat {\mathrm{G}}} (\beta) (x) d \alpha (x)} \\ & {\qquad = \int_ {\mathrm{G}} \overline {{\langle \chi , x \rangle}} \Big (\int_ {\widehat {\mathrm{G}}} \langle \xi , x \rangle d \beta (\xi) \Big) d \alpha (x).} \end{array}
\]

函数 \((x,\xi)\mapsto\overline{\langle\chi,x\rangle}\langle\xi,x\rangle\) 连续且有界，因而在 \(G\times\widehat{G}\) 上关于测度 \(\alpha\otimes\beta\) 可积。由 Lebesgue-Fubini 定理（INT，V，§8，第 4 节，定理 1, a)），我们得到

\[
\begin{array}{r l} & {(\mathcal {F} _ {\mathrm{G}} (\overline {{\mathcal {F}}} _ {\widehat {\mathrm{G}}} (\beta) \cdot \alpha)) (\chi) = \int_ {\widehat {\mathrm{G}}} \Bigl (\int_ {\mathrm{G}} \overline {{\langle \chi \xi^ {- 1} , x \rangle}} d \alpha (x) \Bigr) d \beta (\xi)} \\ & {\qquad = \int_ {\widehat {\mathrm{G}}} \mathcal {F} _ {\mathrm{G}} (\alpha) (\chi \xi^ {- 1}) d \beta (\xi) = (\beta * \mathcal {F} _ {\mathrm{G}} (\alpha)) (\chi),} \end{array}
\]

此即所要证明的。

推论.　Fourier 变换 \(\mathcal{F}_{\mathrm{G}}\) 在 \(\mathcal{M}^{1}(\mathrm{G})\) 上是单射。

事实上，若 \(\alpha \in \mathcal{M}^1 (\mathrm{G})\) 满足 \(\mathcal{F}_{\mathrm{G}}(\alpha) = 0\)，则由 (31) 推出，\(\alpha (\mathcal{F}_{\widehat{\mathrm{G}}} (f)) = 0\) 对每个 \(f\in \mathrm{L}^{1}(\widehat{\mathrm{G}})\) 都成立。由于 \(\mathrm{L}^1 (\widehat{\mathrm{G}})\) 在 Fourier 变换下的像在 \(\mathcal{C}_0(\mathrm{G})\) 中稠密（II，第 209 页，命题 5 的推论），因此有 \(\alpha = 0\)。

在 G 与 \(\widehat{\mathbf{G}}\) 上，除 \(\mathrm{L}^2 (\mathrm{G})\) 与 \(\mathrm{L}^2 (\widehat{\mathrm{G}})\) 之外，还存在使 \(\mathcal{F}\) 与 \(\overline{\mathcal{F}}\) 互为逆同构的其他函数空间。下面的定理给出一个例子。我们用 B(G) 表示 \(\mathrm{L}^1 (\mathrm{G})\) 中由元素 \(f\in \mathrm{L}^1 (\mathrm{G})\)（它们满足 \(\mathcal{F}_{\mathrm{G}}(f)\in \mathrm{L}^1 (\widehat{\mathrm{G}})\)）所组成的向量子空间。它是 \(\mathrm{L}^1 (\mathrm{G})\) 的一个子代数。事实上，设 \(f\) 与 \(g\) 属于 B(G)。我们有 \(f*g\in \mathrm{L}^1 (\mathrm{G})\)，且 \(\mathcal{F}_{\mathrm{G}}(f*g) = \mathcal{F}_{\mathrm{G}}(f)\mathcal{F}_{\mathrm{G}}(g)\in \mathrm{L}^1 (\widehat{\mathrm{G}})\)，因为 \(\mathcal{F}_{\mathrm{G}}(f)\in \mathrm{L}^1 (\widehat{\mathrm{G}})\) 而 \(\mathcal{F}_{\mathrm{G}}(g)\in \mathcal{C}_0(\widehat{\mathrm{G}})\)。

定理 3.　Fourier 变换在 B(G) 上的限制诱导一个从 B(G) 到 B( \(\widehat{G}\) ) 上的向量空间同构，其逆由 Fourier 余变换在 B( \(\widehat{G}\) ) 上的限制所诱导。

设 \(f \in \mathrm{B}(\mathrm{G})\)。记 \(g = \mathcal{F}_{\mathrm{G}}(f)\)。我们有 \(g \in \mathrm{L}^1(\widehat{\mathrm{G}}) \cap \mathcal{C}_0(\widehat{\mathrm{G}}) \subset \mathrm{L}^1(\widehat{\mathrm{G}}) \cap \mathrm{L}^2(\widehat{\mathrm{G}})\)。令 \(f_1 = \overline{\mathcal{F}}_{\widehat{\mathrm{G}}} (g) \in \mathrm{L}^2(\mathrm{G})\)。对每个具有紧支集的连续函数 \(h \in \mathcal{K}(\widehat{\mathrm{G}})\)，有 \(h \in \mathrm{L}^1(\widehat{\mathrm{G}}) \cap \mathrm{L}^2(\widehat{\mathrm{G}})\)，且

\[
\begin{array}{r l} & {\int_ {\mathrm{G}} f _ {1} (x) \mathcal {F} _ {\widehat {\mathrm{G}}} (h) (x) d x = \int_ {\mathrm{G}} \overline {{\mathcal {F}}} _ {\widehat {\mathrm{G}}} (g) (x) \overline {{\overline {{\mathcal {F}}} _ {\widehat {\mathrm{G}}} (\overline {{h}}) (x)}} d x} \\ & {\qquad = \int_ {\widehat {\mathrm{G}}} g (\chi) h (\chi) d \widehat {x} (\chi)} \\ & {\qquad = \int_ {\widehat {\mathrm{G}}} \mathcal {F} _ {\mathrm{G}} (f) (\chi) h (\chi) d \widehat {x} (\chi) = \int_ {\mathrm{G}} f (x) \mathcal {F} _ {\widehat {\mathrm{G}}} (h) (x) d x} \end{array}
\]

这里用到了 Plancherel 定理与公式 (31)。由于 \(\mathcal{K}(\widehat{\mathrm{G}})\) 在 \(\mathrm{L}^2 (\widehat{\mathrm{G}})\) 中稠密，它在 Fourier 变换下的像在 \(\mathrm{L}^2 (\mathrm{G})\) 中稠密。因此有 \(f_{1} = f\)（在 \(\mathrm{L}^{1}(\mathrm{G})\) 中）；这表明 \(g\in \mathrm{B}(\widehat{\mathrm{G}})\)。

公式 \(f_{1}=f\) 表明复合 \(\overline{F}_{\widehat{G}}\circ F_{G}\) 在 B(G) 上的限制是 B(G) 的恒等映射。交换 G 与 \(\widehat{G}\) 的角色，即见 \(F_{G}\circ\overline{F}_{\widehat{G}}\) 是 B( \(\widehat{G}\) ) 的恒等映射，这就完成了定理的证明。

推论 1.　设 \(f \in \mathrm{L}^{1}(\mathrm{G})\)。则 \(f \in \mathrm{B}(\mathrm{G})\) 当且仅当 \(f\) 属于 Fourier 变换 \(\mathcal{F}_{\widehat{\mathrm{G}}}\) 在 \(\mathrm{L}^{1}(\widehat{\mathrm{G}})\) 上的像。特别地，我们有 \(\mathrm{A}(\mathrm{G}) \subset \mathrm{B}(\mathrm{G})\)。

定理 3 证明：若 \(f \in \mathrm{B}(\mathrm{G})\)，则 \(f = \mathcal{F}_{\widehat{\mathrm{G}}}(\overline{\mathcal{F}}_{\mathrm{G}}(f))\)，其中 \(\overline{\mathcal{F}}_{\mathrm{G}}(f)\) 属于 \(\mathrm{L}^1(\widehat{\mathrm{G}})\)。反之，若 \(f = \mathcal{F}_{\widehat{\mathrm{G}}}(g)\)，其中 \(g \in \mathrm{L}^1(\widehat{\mathrm{G}})\)，则由该定理有 \(g \in \mathrm{B}(\widehat{\mathrm{G}})\)，从而 \(f \in \mathrm{B}(\mathrm{G})\)。最后一个断言则由 II，第 217 页，命题 11 得到。

推论 2.　向量空间 B(G) 对于乘法与对于卷积都是一个代数。Fourier 变换把 B(G) 与 B( \(\widehat{G}\) ) 中的卷积与乘法互换。

我们已经看到 B(G) 是 L \(^{1}\) (G) 的一个子代数。另一方面，若 \(f\) 与 \(g\) 属于 B(G)，则 \(fg \in \text{L}^{1}(\text{G})\)，因为 \(f \in \text{L}^{1}(\text{G})\)，而 \(g\) 属于 Fourier 变换在 L \(^{1}\) ( \(\widehat{\text{G}}\) ) 上的像（推论 1）。由于存在 \(f_{1}\) 与 \(g_{1}\)，它们属于 L \(^{1}\) ( \(\widehat{\text{G}}\) )，使得 \(f = \mathcal{F}_{\widehat{\text{G}}} (f_{1})\) 且 \(g = \mathcal{F}_{\widehat{\text{G}}} (g_{1})\)（同前），我们有 \(fg = \mathcal{F}_{\widehat{\text{G}}} (f_{1} * g_{1})\)，于是再由上一个推论得 \(fg \in \text{B(G)}\)。

命题 14.　设 \(f\) 与 \(g\) 属于 \(\mathrm{L}^2 (\mathrm{G})\)。则 \(\mathcal{F}_{\mathrm{G}}(fg) = \mathcal{F}_{\mathrm{G}}(f)*\mathcal{F}_{\mathrm{G}}(g)\)。

若 \(f\) 与 \(g\) 属于 B(G)，则等式成立（推论 2）；特别地，由于 A(G) ⊂ B(G)（推论 1），若 \(f\) 与 \(g\) 属于 A(G)，等式也成立。由于 A(G) 在 L²(G) 中稠密（II，第 212 页，推论 1），只需证明等式的两边都是 \((f,g)\in\mathrm{L}^{2}(\mathrm{G})\times\mathrm{L}^{2}(\mathrm{G})\) 的、取值于 \(\mathcal{C}_{0}(\widehat{\mathrm{G}})\) 的连续函数即可。而映射 \((f,g)\mapsto\mathcal{F}_{\mathrm{G}}(fg)\) 是把连续映射 \((f,g)\mapsto fg\)（由 \(\mathrm{L}^{2}(\mathrm{G})\times\mathrm{L}^{2}(\mathrm{G})\) 到 \(\mathrm{L}^{1}(\mathrm{G})\) 中）与 Fourier 变换 \(\mathcal{F}_{\mathrm{G}}\)（由 \(\mathrm{L}^{1}(\mathrm{G})\) 到 \(\mathcal{C}_{0}(\widehat{\mathrm{G}})\) 中，它也是连续的）复合而成的。同样，映射 \((f,g)\mapsto\mathcal{F}_{\mathrm{G}}(f)*\mathcal{F}_{\mathrm{G}}(g)\) 是把连续映射 \((f,g)\mapsto(\mathcal{F}_{\mathrm{G}}(f),\mathcal{F}_{\mathrm{G}}(g))\)（由 \(\mathrm{L}^{2}(\mathrm{G})\times\mathrm{L}^{2}(\mathrm{G})\) 到 \(\mathrm{L}^{2}(\widehat{\mathrm{G}})\times\mathrm{L}^{2}(\widehat{\mathrm{G}})\) 中）与映射 \((h_{1},h_{2})\mapsto h_{1}*h_{2}\)（由 \(\mathrm{L}^{2}(\widehat{\mathrm{G}})\times\mathrm{L}^{2}(\widehat{\mathrm{G}})\) 到 \(\mathcal{C}_{0}(\widehat{\mathrm{G}})\) 中）（INT，VIII，§4，第 5 节，命题 15）复合而成的。

注.　关于 Fourier 变换在其上为同构的函数空间的其他例子（就特殊的群 G 而言），见第 9 节以及 II，第 270 页习题 22 与 II，第 275 页习题 31。

